\chapter{The Strategy Engine}\label{ch:ll:the-strategy-engine}

Run the same quoting strategy twice on the same recorded day and compare what it sent. If one component stamps its orders
with the wall clock, the two runs never agree, a hundred reruns give a hundred different outputs. If another keeps its quotes
in a hash set whose seed is drawn at start-up and walks them in the set's order, 93 reruns out of a hundred differ. If two
threads feed its queue, 49 do. Remove all three and the hundred runs agree to the last bit, and from then on every bug seen in
production can be replayed on a desk. This chapter builds the component that runs the firm's strategies: a single thread
that processes market events, order reports, timers and parameter changes one at a time, each to completion, in an order
fixed by rules rather than by the operating system, with time injected rather than read, and with nothing on its hot path that
can wait, allocate or vary.

\section{The event loop}

\begin{definition}[Event loop, run-to-completion processing]\label{def:ll:the-strategy-engine:loop}
An \emph{event loop}\index{event loop} is a thread that repeatedly takes the next event from its inputs and hands it to the code
that handles it. \emph{Run-to-completion processing}\index{run-to-completion processing} handles each event entirely, including
everything it causes, before the next event is taken: no handler is interrupted by another, so a handler never sees a state
that another has half changed.
\end{definition}

The engine's inputs are the rings of chapter~12: market events from the \omterm{def:ll:the-feed-handler:handler}{feed handler} (through the \omterm{def:ll:the-order-book-builder:builder}{book builder} of
chapter~19), reports from the order gateway (chapter~21), and control messages; its own timers come from a timing wheel. Its
outputs go to the gateway's ring and to the log. One thread owns all the strategy's state, so nothing needs a lock, and the
order in which events are handled is decided by one rule: before an event with time $t$ is handled, every timer due at or
before $t$ fires, in order of time and then of identifier; parameter changes effective at or before $t$ are applied; then the
event. \Cref{fig:ll:the-strategy-engine:loop} draws it, and \cref{lst:ll:the-strategy-engine:run} is the whole loop.

\begin{omfigure}
\begin{tikzpicture}[font=\footnotesize,box/.style={draw,rounded corners=2pt,minimum height=0.7cm,align=center,inner sep=3pt},
  >=Stealth]
\node[box,fill=omDef!15] (mkt) at (0,1.2) {market events\\ (book builder)};
\node[box,fill=omDef!15] (rep) at (0,0) {order reports\\ (gateway)};
\node[box,fill=omDef!15] (ctl) at (0,-1.2) {parameter\\ snapshots};
\node[box,fill=omMeth!15] (wheel) at (3.6,-2.5) {timing wheel};
\node[box,fill=omThm!12,minimum width=4.2cm,minimum height=2.4cm] (loop) at (4.6,0) {};
\node[align=center] at (4.6,0.7) {\textbf{one thread}};
\node[align=left,font=\scriptsize,anchor=north] at (4.6,0.35) {1. timers due by $t$\\ 2. parameters effective by $t$\\ 3. the event at $t$};
\node[box,fill=omDef!15] (gw) at (9.0,0.6) {orders\\ (gateway ring)};
\node[box,fill=gray!10] (log) at (9.0,-0.6) {log};
\draw[->] (mkt) -- (loop);
\draw[->] (rep) -- (loop);
\draw[->] (ctl) -- (loop);
\draw[->] (wheel.north) -- (wheel.north |- loop.south);
\draw[->,dashed] ([xshift=0.9cm]loop.south) |- node[pos=0.75,below,font=\scriptsize]{schedule} (wheel.east);
\draw[->] (loop) -- (gw);
\draw[->] (loop) -- (log);
\node[font=\scriptsize,align=left,anchor=west] at (7.4,-2.0) {clock: injected\\ (simulated or real)};
\end{tikzpicture}

\omcaption{The strategy engine. Every input reaches one thread, which handles each event to completion in an order fixed by
rule: the timers due by the event's time, then the parameters effective by it, then the event. Time is whatever clock the
engine was given: a recorded day's in replay, the machine's live.}\label{fig:ll:the-strategy-engine:loop}
\end{omfigure}

\omcode{code/firm/stratengine/cpp/firm_stratengine.hpp}{128}{136}{The event loop: timers first, then parameters, then the
event, each handled to completion.}\label{lst:ll:the-strategy-engine:run}

The design is old and deliberately plain. The retail exchange whose architecture Martin Fowler described in 2011 ran its
business logic on one thread, in memory, with the state ``entirely derivable by processing the input events'' kept in a
durable store; and because ``the business logic is deterministic'', its outputs did not even need to be stored. A strategy
engine gains the same two things: speed, because one thread with its data in its own caches does not coordinate with
anyone, and replay, because the same inputs in the same order produce the same outputs.

\section{Single-threaded determinism}

A program is reproducible bit for bit (the bitwise reproducibility of One Quant Book~4, chapter~25) only as a whole, and any
one of its parts can break it. \Cref{fig:ll:the-strategy-engine:reruns} counts the damage on the build's small quoting engine, run a hundred times as a new
process on the same 6\,000 recorded events, with one source of nondeterminism switched on at a time.

\begin{itemize}
\item \emph{A read of the wall clock}, here to stamp each action and to date its quotes: every run reads different times, so a
hundred runs give a hundred outputs. Any decision that compares the real time with something (a quote's age, a timeout) makes
the difference reach the orders themselves.
\item \emph{Iteration over a hash container with a per-process seed}: the engine keeps its ten quotes in a hash set whose hash
is seeded at start-up and requotes them in the set's order, so the identifiers and the order of the actions follow the seed.
Rust's standard hash map does exactly this by design (its algorithm ``is randomly seeded''); 93 of the hundred runs
differ.
\item \emph{Two threads feeding one queue}, each carrying half of the events: the order in which the consumer sees them depends
on scheduling, the book passes through states that never existed, and 49 of the hundred runs differ.
\end{itemize}

Switched off, all three give one output in a hundred runs; so does the firm's engine. The rules that achieve it are few. Time
comes from the events, through an \omterm{def:ll:the-strategy-engine:wheel}{injected clock} (next section). Containers that are iterated have a defined order: arrays,
ordered maps, or hash maps whose order does not decide anything. One thread owns the state; other threads only move data
through rings, and the merge of several inputs follows a rule (by time, then by source, then by sequence), not the arrival
order. Floating-point sums are done in a fixed order (chapter~8), and the example strategy stays in integers, which is also why
the C++ and Rust engines of the build produce the same output hash.

\begin{omfigure}
\begin{tikzpicture}
\begin{axis}[width=11cm,height=5.4cm,ybar,bar width=16pt,ymin=0,ymax=115,
  symbolic x coords={firm,none,wall_clock,hash_order,two_threads,all},xtick=data,
  xticklabels={{firm's\\ engine},{no\\ defect},{wall\\ clock},{seeded\\ hash order},{two\\ threads},{all\\ three}},
  xticklabel style={align=center,font=\footnotesize},enlarge x limits=0.1,ylabel={distinct outputs in 100 runs},
  tick label style={font=\small},label style={font=\small},grid=major,grid style={gray!25},
  nodes near coords,every node near coord/.append style={font=\scriptsize},table/col sep=comma]
\addplot[fill=omThm!45,draw=omThm] table[x=key,y=distinct]{figdata/low-latency/20-the-strategy-engine/measured_nondet.csv};
\end{axis}
\end{tikzpicture}

\omcaption{Distinct output hashes in a hundred reruns (each a new process) of the same 6\,000 recorded events, for the firm's
engine and for the small engine of \texttt{ll\_nondet.hpp} with each source of nondeterminism switched on. Measured on a laptop
(Intel Core Ultra 7 155H) under WSL2. Data: \texttt{bench\_engine.py}.}\label{fig:ll:the-strategy-engine:reruns}
\end{omfigure}

\section{Timers and the injected clock}

\begin{definition}[Timer wheel, injected clock]\label{def:ll:the-strategy-engine:wheel}
A \emph{timer wheel}\index{timer wheel} is a circular array of slots, each covering an interval of time (its granularity): a
timer is filed in the slot of its expiry time modulo the wheel's span, and advancing the wheel visits the slots of the time that
has passed and fires the timers due in them; scheduling costs a constant time whatever the number of timers. An \emph{injected
clock}\index{injected clock} is a clock passed to the code that needs the time instead of being read by it: the same code runs
on the simulated clock of a recorded day (One Quant Book~7, chapter~17) or on the machine's real clock.
\end{definition}

A strategy is full of timers: an order that must be acknowledged within a delay, a quote that must be pulled if the market data
stops, a refresh every few seconds, a \omterm{def:ll:protocols-i-fix:session}{heartbeat}. A priority queue would do, at a logarithmic cost per timer; the hashed timing
wheel of Varghese and Lauck (1987) does it in constant time, which matters when every order schedules and cancels timers.
The build's wheel has 1\,024 slots of \qty{100}{\micro\second}: a timer goes to slot $\lfloor t/100\,\mu\text{s}\rfloor \bmod 1\,024$,
and a timer further away than the wheel's span of about \qty{0.1}{\second} simply stays in its slot until its time comes round.
Two details make it deterministic: the timers due in a slot fire sorted by time and then by identifier, never in insertion
order; and a timer scheduled while another fires, if due in the same slot, fires in the same pass.

\begin{omfigure}
\begin{tikzpicture}[font=\scriptsize,>=Stealth]
\foreach \i in {0,...,15} {
  \pgfmathsetmacro{\a}{90 - 22.5*\i}
  \pgfmathtruncatemacro{\busy}{(\i==2)||(\i==3)||(\i==7)||(\i==12)}
  \ifnum\busy=1 \draw[fill=omMeth!25,draw=omMeth] (\a+11.25:1.4) arc (\a+11.25:\a-11.25:1.4) -- (\a-11.25:2.3) arc (\a-11.25:\a+11.25:2.3) -- cycle;
  \else \draw[fill=gray!6,draw=gray!50] (\a+11.25:1.4) arc (\a+11.25:\a-11.25:1.4) -- (\a-11.25:2.3) arc (\a-11.25:\a+11.25:2.3) -- cycle; \fi
  \node at (\a:1.85) {\i};
}
\draw[->,very thick,omThm] (0,0) -- (45:1.35);
\node[align=center] at (0,-0.35) {now: slot 2};
\node[anchor=west,align=left] at (2.9,1.4) {slot $= \lfloor t / g \rfloor \bmod 16$ ($g$: granularity)};
\node[anchor=west,align=left] at (2.9,0.7) {slot 3: ack of quote 41 at $t_0 + 50\,\mu$s;\\ stale check at $t_0 + 20$ ms};
\node[anchor=west,align=left] at (2.9,-0.2) {advancing to $t$: visit slots 2, 3, \dots\ up to $t$'s,\\ fire the timers due by $t$ in (time, id) order;};
\node[anchor=west,align=left] at (2.9,-1.1) {a timer more than a revolution away\\ (the stale check) waits in its slot};
\end{tikzpicture}

\omcaption{A hashed timing wheel of 16 slots (the build's has 1\,024 of \qty{100}{\micro\second}). Scheduling puts a timer in
the slot of its expiry time; advancing visits the slots passed and fires only the timers due, sorted; a timer further than one
revolution stays in its slot until its time comes round.}\label{fig:ll:the-strategy-engine:wheel}
\end{omfigure}

\omcode{code/firm/stratengine/cpp/firm_stratengine.hpp}{82}{101}{Advancing the wheel: visit the slots passed, fire the timers
due in time and identifier order.}\label{lst:ll:the-strategy-engine:wheel}

The wheel knows no clock. It is advanced to the time of each event, and the engine's code asks an \omterm{def:ll:the-strategy-engine:wheel}{injected clock} for the time:
the simulated clock returns the event's time, the real clock reads \texttt{CLOCK\_MONOTONIC} (chapter~5) and maps it onto the
day. Replaying a day with the simulated clock gives, event for event, what the live engine would have done with the same inputs
at the same times; the real clock is used live, and only there.

\omcode{code/firm/stratengine/cpp/firm_stratengine.hpp}{46}{58}{The two clocks the engine can be given.}\label{lst:ll:the-strategy-engine:clock}

\section{Configuration and parameters at run time}

\begin{definition}[Parameter snapshot]\label{def:ll:the-strategy-engine:snapshot}
A \emph{parameter snapshot} is a complete, versioned set of a strategy's parameters that takes effect at a stated time, between
two events and never during one; every action the strategy takes records the version of the parameters it was taken under.
\end{definition}

Traders change parameters during the day: a wider spread before a number, a smaller size in a thin market. Changing a field of
a live structure from another thread while the engine reads it is a \omterm{def:ll:rust-for-low-latency:race}{data race} (chapter~11), and even a correct atomic update of
one field leaves a moment where half the parameters are new. The build applies whole snapshots, in the loop, between events:
the control path publishes a snapshot with an effective time on the engine's control ring, and the loop installs it before the
first event at or after that time. The version stamped on every action says which parameters produced it, and in replay the
snapshots are inputs like any other event: the build's test applies one half-way through the day and checks that every action
before it is unchanged and every later one carries the new version.

\section{Hot-path discipline}

The hot path of chapter~1 runs through this thread, and it is disciplined by rules the earlier chapters have measured. No
allocation per event (chapter~6): the build's test counts heap allocations while the engine replays its fixture and finds none,
because the action log, the wheel's slots and the book are sized at start-up. No system call: the simulated clock is a field,
and even the real clock is read through the vDSO (chapter~5). No lock and no waiting: inputs and outputs are rings, and a full
output ring is an error to report, not a reason to wait. No formatting of text: the log receives binary records (chapter~23). No
unbounded work in a handler: the wheel's work is bounded by the slots passed, the book's by its structure.

The result on the laptop: the engine takes about \qty{53}{\nano\second} per event at the median on the large-tick stream of
chapter~18, book update included, and about \qty{57}{\nano\second} on the small-tick one where it requotes at almost every
event; the p99 is about \qty{130}{\nano\second} and \qty{200}{\nano\second}. A decision in tens of nanoseconds leaves the
\omterm{def:ll:where-latency-comes-from:budget}{latency budget} to the network and the gateway.

\section{Tutorial: replay a day, then break it}\label{tut:ll:the-strategy-engine:tut}

\begin{tutorial}
\textbf{Goal.} Replay a recorded input through the engine to the same output every time, then switch on each source of
nondeterminism and count the outcomes. \textbf{End state:} \cref{fig:ll:the-strategy-engine:reruns}, the engine's time per
event, and green tests in C++ and Rust.
\begin{tutsteps}
\item \textbf{Replay.} The C++ test runs the engine twice on the \omterm{def:ll:the-order-book-builder:builder}{book builder}'s small-tick fixture and compares the output
hashes, the actions and the profit, then checks them against \texttt{data/expected.txt}; the Rust engine, written against the
same rules, reproduces the same hash.
\item \textbf{Change the parameters half-way} with a snapshot and check that only the actions after it change.
\item \textbf{Count the zero:} the test counts heap allocations during the replay.
\item \textbf{Break it} with \texttt{python bench\_engine.py}: a hundred new processes per configuration of \texttt{ll\_nondet},
each printing its output hash, then the distinct hashes counted.
\end{tutsteps}
\textbf{What to change next.} Replace the small engine's seeded hash set by one with a fixed seed and count again; feed the two
threads' events through a merge by timestamp instead of a shared queue.
\end{tutorial}

\section{Build: the strategy engine}\label{bld:ll:the-strategy-engine:stratengine}

\begin{build}
\bfield{Purpose} The thread on which the firm's strategies run: it consumes the \omterm{def:ll:the-order-book-builder:builder}{book builder}'s view (chapter~19), sends orders
through the gateway (chapter~21) past the risk gate (chapter~22), and logs everything it does (chapter~23).
\bfield{Interface} C++20 \texttt{firm::strat}: \texttt{Engine<Clock>(tick, Params, clock)} with \texttt{run(events, snapshots)},
\texttt{actions}, \texttt{hash}, \texttt{position}, \texttt{cash}, \texttt{pnl()}; \texttt{TimerWheel(granularity, slots)} with
\texttt{schedule} and \texttt{advance}; \texttt{SimClock}, \texttt{RealClock}; \texttt{Params}, \texttt{Snapshot},
\texttt{Action}. Rust \texttt{firm\_stratengine}: \texttt{Engine}, \texttt{TimerWheel}, \texttt{Params}, \texttt{Action},
\texttt{read\_events}.
\bfield{Rules} One thread owns the state; each event is handled to completion; timers due by an event fire before it, in (time,
identifier) order; parameter snapshots apply between events; time comes from the \omterm{def:ll:the-strategy-engine:wheel}{injected clock}; the example strategy computes
in integers; no allocation per event.
\bfield{Acceptance tests} \texttt{code/firm/stratengine/}: two replays with identical actions, hash and profit, equal to the
committed hash; the Rust engine equal to it; a snapshot changing only what follows it; zero allocations during the replay; the
wheel firing in order across slots and revolutions.
\bfield{Stretch} A hierarchical wheel for long timers; order reports from the simulator of One Quant Book~10 in place of the venue
stub; several strategies on one engine with a fixed scheduling order.
\end{build}

\begin{omsources}
\item G.~Varghese and T.~Lauck, ``Hashed and hierarchical timing wheels: data structures for the efficient implementation of a
timer facility'', \emph{Proceedings of the 11th ACM Symposium on Operating Systems Principles}, 1987.
\item M.~Fowler, \emph{The LMAX Architecture}, 2011.
\item The Rust standard library, \texttt{std::collections::HashMap}.
\end{omsources}

\section{Exercises}

\begin{exercise}[$\star$]\label{exo:ll:the-strategy-engine:1}
A wheel has 1\,024 slots of \qty{100}{\micro\second}. In which slot does a timer due at 34\,200.012345 seconds after midnight go,
and how many revolutions away is a timer due in 5 seconds?
\end{exercise}

\begin{exercise}[$\star$]\label{exo:ll:the-strategy-engine:2}
Three timers are due in the same slot at 120, 121 and 120 microseconds with identifiers 7, 3 and 9. In what order do they fire?
\end{exercise}

\begin{exercise}[$\star$]\label{exo:ll:the-strategy-engine:3}
Why is the time of an event the right ``now'' for a strategy in replay, and what must be true of the live system for replay to
match it?
\end{exercise}

\begin{exercise}[$\star\star$]\label{exo:ll:the-strategy-engine:4}
If an outcome depends on which of $k$ equally likely orders a hash set produces, how many distinct outcomes should a hundred
reruns show on average? What does the measured 93 suggest about $k$?
\end{exercise}

\begin{exercise}[$\star\star$]\label{exo:ll:the-strategy-engine:5}
A trader changes the spread from two ticks to four at 10:30:00.000. An event arrives at 10:29:59.999 and the next at
10:30:00.004. Under which parameters is each handled, and what does the log show?
\end{exercise}

\begin{exercise}[$\star\star$]\label{exo:ll:the-strategy-engine:6}
The engine takes \qty{53}{\nano\second} per event at the median. At the open, events arrive at 2 million a second for 50
milliseconds. Does it keep up, and what is its utilisation?
\end{exercise}

\begin{exercise}[$\star\star\star$]\label{exo:ll:the-strategy-engine:7}
\emph{Coding.} Add a refresh timer to the engine: every quote older than \qty{5}{\second} is re-sent at the same price. Check
that two replays still agree and that the Rust engine, changed the same way, still produces the C++ hash.
\end{exercise}

\begin{exercise}[$\star\star\star$]\label{exo:ll:the-strategy-engine:8}
\emph{Find the flaw.} ``Our engine is deterministic: it is single-threaded and uses no randomness. It stamps each order with
\texttt{std::chrono::system\_clock::now()} for the audit trail.''
\end{exercise}

\section{Problem: Two Runs, Two Profits}

\begin{problem}[{Weekend problem --- counting the ways a program can disagree with itself}]\label{pb:ll:the-strategy-engine:1}
A research team replays the same recorded day through its quoting engine twice and gets two profits. Use the measurements of
\cref{fig:ll:the-strategy-engine:reruns} (\texttt{measured\_nondet.csv}).

\medskip\noindent\textbf{Part I --- The sources.}
\begin{enumerate}
\item List the three sources measured and the number of distinct outputs each gave in a hundred runs.
\item Why does the wall clock give a hundred?
\item Why does the seeded hash set give fewer than a hundred?
\item Why does the two-thread feed give fewer still, and what would make it give more?
\end{enumerate}

\medskip\noindent\textbf{Part II --- Finding them.}
\begin{enumerate}[resume]
\item How would you find the first action at which two runs differ?
\item What does each source leave as a clue at that point?
\item Why must each rerun be a new process to see the seeded hash's effect?
\item Which of the three sources would a single run inside a debugger never show?
\end{enumerate}

\medskip\noindent\textbf{Part III --- The fixes.}
\begin{enumerate}[resume]
\item Fix the wall clock.
\item Fix the hash order without giving up hash maps.
\item Fix the two-thread feed.
\item After the fixes, how many distinct outputs should a hundred reruns give, and how many did the firm's engine give?
\end{enumerate}

\medskip\noindent\textbf{Part IV --- The verdict.}
\begin{enumerate}[resume]
\item State the \emph{named result}: the number of distinct outcomes in a hundred reruns with each source switched on, and after
the fixes.
\item Why is ``the profit differs by a little'' not a reason to accept nondeterminism in a backtest?
\item What else can break bitwise reproducibility between two machines, even with all three fixes?
\item Why does integer arithmetic in the strategy make the C++ and Rust engines agree?
\item What should the continuous-integration pipeline run to keep determinism?
\item What does determinism give the production support team?
\item When is a nondeterministic component acceptable in a trading system?
\item In one sentence: what makes an engine replayable?
\end{enumerate}
\end{problem}

\section{Interview questions}

\begin{interviewq}[$\star$ \iqroles{developer}]\label{iq:ll:the-strategy-engine:1}
Why run a strategy on a single thread? What does it cost?
\end{interviewq}

\begin{interviewq}[$\star\star$ \iqroles{developer}]\label{iq:ll:the-strategy-engine:2}
How would you implement timers for thousands of orders, each with its own timeouts?
\end{interviewq}

\begin{interviewq}[$\star\star$ \iqroles{developer}]\label{iq:ll:the-strategy-engine:3}
Your backtest and your production run disagree on the same day. Where do you look?
\end{interviewq}

\begin{interviewq}[$\star\star$ \iqroles{developer}]\label{iq:ll:the-strategy-engine:4}
How do you change a strategy's parameters while it runs, safely?
\end{interviewq}

\begin{interviewq}[$\star\star$ \iqroles{developer}]\label{iq:ll:the-strategy-engine:5}
What must never happen on the hot path of a strategy engine, and how do you check?
\end{interviewq}

\begin{interviewq}[$\star\star\star$ \iqroles{developer, researcher}]\label{iq:ll:the-strategy-engine:6}
Design a strategy engine that is both live and replayable bit for bit. What are the rules, and how are they enforced?
\end{interviewq}
